\section{Orbit closures and compact targets}
\label{sec:closure}

The direct image can be nonclosed.  Its closure may contain new recurrent or
periodic points, so exactness on $\pi(X)$ alone cannot control them.  We now
freeze the additional inheritance condition needed at the boundary.

Let $Y_0=\cl{\pi(X)}$ in a topological space $Y$, assume $S:Y_0\to Y_0$ is
continuous, and let $C$ be a topological clock space.  For $m\ge1$ define
\begin{equation}
P_m=\{(a_k,a_{k+m}):k\ge0\}\subset C\times C,
\qquad
\diag_C=\{(c,c):c\in C\}.
\label{eq:lag-pairs}
\end{equation}

\begin{theorem}[Continuous closure-clock obstruction]
\label{thm:closure}
Suppose $d:Y_0\to C$ is continuous and satisfies
$d(\pi(x))=a_{\ell(x)}$.  If
\begin{equation}
\cl{P_m}\cap\diag_C=\varnothing,
\label{eq:diagonal-separation}
\end{equation}
then $Y_0\cap\Per_m(S)=\varnothing$.
\end{theorem}

\begin{proof}
Define the continuous map
\[
F_m:Y_0\to C\times C,
\qquad
F_m(y)=(d(y),d(S^my)).
\]
Equivariance and exact decoding imply
$F_m(\pi(X))=P_m$.  Since $\pi(X)$ is dense in $Y_0$, continuity gives
$F_m(Y_0)\subseteq\cl{P_m}$.  If $S^my=y$, then
$F_m(y)=(d(y),d(y))\in\diag_C$, contradicting
\eqref{eq:diagonal-separation}.
\end{proof}

No metrizability, first countability, or continuity of $\pi$ is needed in
this proof.  The relevant regularity is continuity of the target dynamics
and of the inherited clock.

\begin{corollary}[Exact wheel clocks]
\label{cor:closure-wheel}
\Cref{thm:closure} applies to each of the following:
\begin{enumerate}[label=(\roman*)]
  \item $C=\N$ with the discrete topology and $a_k=q_{k+1}$;
  \item $C=\R$ with its usual topology and $a_k=q_{k+1}$;
  \item $C=\R$ with its usual topology and $a_k=\log q_{k+1}$.
\end{enumerate}
Consequently a continuous exact decoder of any of these clocks excludes all
periodic points from $Y_0$.
\end{corollary}

\begin{proof}
In the discrete case, $P_m$ is closed and disjoint from the diagonal.  In
the real cases, both coordinates tend to $+\infty$ with $k$, so $P_m$ is
locally finite in $\R^2$ and has no finite accumulation point.  Hence it is
closed in $\R^2$, and pairwise distinctness excludes its intersection with
the diagonal.
\end{proof}

Compactness gives a separate feasibility obstruction, independent of the
periodic-point argument.

\begin{proposition}[Compact-target obstruction]
\label{prop:compact}
Let $Y_0$ be compact and contain $\pi(X)$.  There is no continuous total
exact wheel-clock decoder $d:Y_0\to\N_{\mathrm{disc}}$.  There is likewise
no continuous decoder $d:Y_0\to\R$ of all values $q_{k+1}$ or
$\log q_{k+1}$.
\end{proposition}

\begin{proof}
The continuous image of a compact space is compact.  A compact subset of a
discrete space is finite, while a compact subset of $\R$ is bounded.  Exact
decoding would place the infinite unbounded clock range in that image.
\end{proof}

This proposition uses only compactness of $Y_0$ and continuity of $d$; it
uses neither equivariance, periodicity, nor the diagonal-separation theorem.

Thus a compact finite-alphabet shift cannot continuously retain the complete
absolute clock, regardless of decoder memory.  A compactified or transformed
clock may have a continuous extension, but it must then be checked against
\eqref{eq:diagonal-separation}; continuity alone is not enough.
